% ============================================================
\chapter{Connections to TDA}
\label{ch:tda}
% ============================================================

\begin{intuition}
Topological Data Analysis (TDA)\index{TDA} and geometric deep
learning share a common concern: exploiting the \textbf{structure}
of data. TDA provides topological invariants (Betti numbers,
persistent homology) that complement features learned by GNNs
and improve their \textbf{expressiveness}.
\end{intuition}

% ============================================================
\section{Review of persistent homology}
\label{sec:persistent_homology_review}
% ============================================================

\begin{definition}[Persistent homology]
\index{persistent homology}
Given a filtration of simplicial complexes
$K_0 \subseteq K_1 \subseteq \cdots \subseteq K_m$,
\textbf{persistent homology} tracks the evolution of homology
classes through the filtration. Each class is born at index $b$
and dies at index $d > b$, yielding a point $(b, d)$ in the
\textbf{persistence diagram}.
\end{definition}

\begin{remark}
For a graph $G = (V, E)$, different filtrations can be constructed:
\begin{itemize}
  \item \textbf{By degree}: $f(v) = \mathrm{deg}(v)$ and
        $f(e) = \max(f(u), f(v))$.
  \item \textbf{By centrality}: $f(v) = \mathrm{betweenness}(v)$.
  \item \textbf{By learned features}: $f(v) = g_\theta(h_v)$
        where $g_\theta$ is a neural network.
\end{itemize}
\end{remark}

% ============================================================
\section{Topological features in GNNs}
\label{sec:topo_features}
% ============================================================

\begin{definition}[Topological features for graphs]
\index{topological features}
\textbf{Topological features} of a graph are extracted from its
persistence diagrams:
\begin{itemize}
  \item $H_0$: connected components and their stability.
  \item $H_1$: cycles and their persistence.
\end{itemize}
These features are vectorized (landscapes, persistence images) and
concatenated with GNN features.
\end{definition}

\begin{proposition}[Complementarity]
Topological features capture global graph properties (cycles,
cavities) that local message passing cannot detect in a bounded
number of iterations: detecting a cycle of length $k$ requires
at least $k/2$ message-passing layers.
\end{proposition}

\begin{minted}{python}
import torch
import numpy as np
from gudhi import RipsComplex
from gudhi.representations import PersistenceImage

def topo_features(pos, max_edge=5.0, resolution=10):
    """Compute topological features of a point cloud."""
    rips = RipsComplex(points=pos, max_edge_length=max_edge)
    st = rips.create_simplex_tree(max_dimension=2)
    st.persistence()

    # Diagrams by dimension
    dgm_h0 = st.persistence_intervals_in_dimension(0)
    dgm_h1 = st.persistence_intervals_in_dimension(1)

    # Vectorization via persistence images
    pi = PersistenceImage(
        bandwidth=0.1, resolution=[resolution, resolution]
    )
    feats = []
    for dgm in [dgm_h0, dgm_h1]:
        if len(dgm) > 0:
            dgm = dgm[np.isfinite(dgm[:, 1])]
        if len(dgm) > 0:
            feats.append(pi.fit_transform([dgm])[0])
        else:
            feats.append(np.zeros(resolution ** 2))
    return np.concatenate(feats)
\end{minted}

% ============================================================
\section{Topological message passing}
\label{sec:topo_message_passing}
% ============================================================

\begin{definition}[Message passing on simplicial complexes]
\index{simplicial message passing}
\textbf{Simplicial message passing} (Bodnar et al.\ 2021)
generalizes graph message passing to higher-dimensional simplicial
complexes:
\[
  h_\sigma^{(t+1)} = \phi\!\left(h_\sigma^{(t)},\;
  \bigoplus_{\tau \in \mathcal{B}(\sigma)}
    \psi_{\mathcal{B}}(h_\tau^{(t)}),\;
  \bigoplus_{\tau \in \mathcal{C}(\sigma)}
    \psi_{\mathcal{C}}(h_\tau^{(t)})
  \right)
\]
where $\mathcal{B}(\sigma)$ are the boundary faces and
$\mathcal{C}(\sigma)$ are the coboundary cofaces of simplex $\sigma$.
\end{definition}

\begin{intuition}
In a graph, messages flow along edges (1-simplices). In a simplicial
complex, messages also flow along triangles (2-simplices) and
higher-dimensional structures. This enables capturing
\textbf{higher-order} relationships between nodes.
\end{intuition}

\begin{definition}[Message passing on cell complexes]
\index{CW Network}
\textbf{CW Networks} (Bodnar et al.\ 2021) extend message passing
to cell complexes, which generalize simplicial complexes by allowing
cells of arbitrary shape (not just simplices).
\end{definition}

% ============================================================
\section{The Weisfeiler-Leman hierarchy and topology}
\label{sec:wl_topology}
% ============================================================

\begin{definition}[Weisfeiler-Leman test]
\index{Weisfeiler-Leman}
The \textbf{Weisfeiler-Leman} (WL) test of level $k$ is a
color refinement algorithm that characterizes graphs: two graphs
not distinguished by $k$-WL are called \textbf{$k$-WL equivalent}.
\end{definition}

\begin{theorem}[GNN expressiveness --- Xu et al.\ 2019, Morris et al.\ 2019]
\index{WL expressiveness}
A message-passing GNN is \textbf{at most as powerful} as the
$1$-WL test for distinguishing non-isomorphic graphs.
\end{theorem}

\begin{theorem}[Beyond WL via topology --- Bodnar et al.\ 2021]
\index{simplicial complex}
Message passing on simplicial complexes with topological features
(persistent homology) can distinguish graphs that $k$-WL cannot
distinguish, for any fixed $k$.
\end{theorem}

\begin{example}[Rook graphs]
The Rook graphs $R_{3 \times 3}$ and $R_{4 \times 4}$ are not
distinguished by $3$-WL. However, their $H_1$ persistence diagrams
(with degree filtration) differ: $R_{4 \times 4}$ has more
persistent cycles.
\end{example}

% ============================================================
\section{Expressiveness and topology}
\label{sec:expressiveness}
% ============================================================

\begin{proposition}[Topological enrichment]
By adding persistent homology features to GNN node features, the
resulting expressiveness is \textbf{strictly greater} than that of
the GNN alone for graph classification.
\end{proposition}

\begin{definition}[Joint learned filtration]
\index{learned filtration}
In a \textbf{joint learned filtration}, filtration values are
defined by the GNN itself:
\[
  f_\theta(v) = \mathrm{MLP}_\theta(h_v^{(L)})
\]
where $h_v^{(L)}$ is the final embedding of node $v$ after $L$
message-passing layers. Persistence is computed on this filtration,
and gradients backpropagate through the persistence computation
(cf.\ differentiable persistence).
\end{definition}

\begin{keyformulas}[title=\textbf{GNN + TDA pipeline}]
\[
  X \xrightarrow{\text{GNN}} H^{(L)}
  \xrightarrow{f_\theta} \text{Filtration}
  \xrightarrow{\text{PH}} \mathrm{Dgm}
  \xrightarrow{\text{PersLay}} z_{\mathrm{topo}}
  \xrightarrow{\oplus} \text{Prediction}
\]
\end{keyformulas}

% ============================================================
\section{Exercises}
% ============================================================

\begin{exercise}
For the Petersen graph, compute the $H_0$ and $H_1$ persistence
diagrams with the degree filtration. All nodes have the same degree:
what can you conclude?
\end{exercise}

\begin{exercise}
Show that two isomorphic graphs have the same persistence diagrams
for any filtration defined from structural graph properties.
\end{exercise}

\begin{exercise}[Simplicial message passing]
Write the message-passing equations for a triangle $(v_1, v_2, v_3)$
in a simplicial complex. What messages does the edge $(v_1, v_2)$
receive?
\end{exercise}

\begin{exercise}
Give an example of two non-isomorphic graphs with the same
persistence diagrams for the degree filtration. Propose a filtration
that distinguishes them.
\end{exercise}

\begin{exercise}[Project]
Implement a GNN enriched with topological features on the ZINC
benchmark and compare with a baseline GCN. Measure the improvement
in MAE.
\end{exercise}
